\section{Challenges} \label{sec:main-challenges}

While the field of AI for code has made fruitful progress, cutting-edge AI still struggles with SWE tasks, especially at larger scopes and higher levels of logical complexity. Next, we discuss four high-level challenges in AI for code: data, scale, interaction, and measurement. These four challenges permeate across all of the tasks mentioned in the previous section.

\subsection{Data}

\textbf{Low-Resource Languages and Specialized Libraries}: As we adapt code LLMs to individual codebases, generating correct code in out of distribution (OOD) scenarios becomes crucial. Much of software development in business contexts revolves around proprietary codebases, which is a distribution shift from the open-source code that dominates LLM training data  \citep{ahmed2024studying}. These OOD scenarios include domain-specific languages (DSLs), custom internal libraries, low-resource APIs, and company-specific coding styles and conventions.

In low-resource languages, models have a weaker semantic understanding of the language. Due to the lack of training data in these OOD domains, models may struggle to write common primitives or piece together functionality coherently. On HumanEval, Qwen-2.5 has an accuracy of 83\% in Python but 27\% in D \citep{hui2024qwen2}. In OOD scenarios, LLMs lack awareness of the libraries and functions available for use. In new codebases using custom libraries, many functions appear only a few times, providing limited training data for AI models to learn their usage. This scarcity can lead to overfitting, where models fail to recognize an effective use-case of these functions. Models also frequently hallucinate non-existent functions based on patterns that it infers.

\textbf{Library and API Version Updates}: Continual learning, the idea of training an AI system to take in new information continually, has been a long-standing challenge in AI and NLP  \citep{wu2024continual, wang2024comprehensive}. In software engineering, codebases are continuously changing as new features are supported and awkward design patterns are reworked. While backwards compatibility is often prioritized in software design, it inevitably becomes broken as codebases evolve further. Therefore, programming libraries have version releases, each release supporting and deprecating features in the last version. 

Even identifying which version of a library is being used can be quite difficult, because versioning information can be hidden deeply within a codebase. Sometimes, it can be found in comments or configuration files, but in the worst case, it must be inferred from the library calls being used. To make things worse, some code may be compatible across multiple versions, while other code will cause errors only in specific versions. On the other hand, writing code for a specific version can also be challenging, because it can be difficult for LLMs to implicitly keep track of which constructs and patterns are associated with each version. For example, when asking LLMs to write code in Lean 4, it often uses constructs from Lean 3, an older version that is much more prevalent across GitHub.

\textbf{High Logical Complexity and OOD Domains}: Some programming tasks are challenging for even the best human programmers, requiring approaches with a very high logical complexity. Examples of tasks that fall into this category include superoptimizing programs, discovering attacks for purportedly secure code, writing performant compilers, optimizing GPU kernels \citep{ouyang2024kernelbench}, and writing concurrent programs.

Because they are hard for humans, these tasks are very rarely in the training data of today's language models. They have unique, domain-specific, challenges that making generalizing from existing data difficult. For these problems, language models rely heavily on feedback-driven search algorithms \citep{mankowitz2023faster}, and it can be difficult to navigate the search space effectively. 

\subsection{Scale}

\textbf{Large Scope and Long Contexts}:
At the repository level, the tasks in Sec. \ref{sec:main-tasks-milestones} become significantly more difficult and require many steps. In code generation, user alignment can be an issue because there are many decision points and tradeoffs that can compound. In code refactoring, modifications will touch multiple parts of the codebase, and it can be tricky to keep the repository consistent. In code debugging, functions can be large and bugs can be nested deeply within stacks of function calls.
Another issue with large scopes is large context lengths. Software engineering often requires dealing with very large codebases--for example, Google has repositories with over a billion lines of code \citep{potvin2016google}. As this is far too large for modern-day LLMs, choosing the correct context to include when using LLMs is important.

\textit{Limits of Retrieval-Augmented Generation (RAG)}: Retrieval-based algorithms have been the predominant way to deal with long-context coding issues. First, the retriever retrieves relevant functions. Then, the generator leverages the retrieval to improve generation. While RAG has proven effective in many NLP tasks, the code domain provides new challenges for these methods.

\textit{Retrieval}: In most NLP tasks, the retrieval step can be done relatively well because keywords that are in the query are often keywords that need to be retrieved. Unlike answering NL questions, writing code often requires drawing inspiration from code snippets that may be completely different syntactically. This can include programs with similar semantics, algorithms, or API calls, all of which potentially have very little in common when it comes to syntax. Previous work such as \citet{ma2024unveiling, utpala2023language} have found that code in the same language is, on average, far closer in embedding space compared to semantically equivalent code.

\textit{Generation}: In NLP tasks, the generation step often is a straightforward application of the retrieved information. In code, we are dealing with code reuse: piecing together relevant snippets of code in a precise and productive way to fit the current context. Each piece of retrieved content provides different information. This can include information about the language's syntax, documentation about the API, clues about the algorithm to be written, or examples of similar functionality being written. 

\textbf{Long-Horizon Code Planning}
When working on large projects, engineers often make complex decisions about how to design and structure the code to best support the various functionalities that will eventually be needed. To build a long-lasting software system, an engineer must know the potential paths that the system's evolution might take. This requires domain expertise and experience in how different code structures require different forms of extension. 

\textit{Designing Good Abstractions}:
One instance of long-horizon code planning is choosing the right abstractions from the outset. An API designed with good abstractions will allow new features to be implemented seamlessly with minimal user overhead, while an API designed with poor abstractions may lead to excessive code duplication, refactoring, or even debugging. An example of this is \textit{library learing}: finding the right APIs and libraries that can provide useful abstractions, leading to more reuse and intuitive interfaces. \citep{ellis2021dreamcoder,10.5555/3692070.3693967,10.1145/3571234}. While the traditional library learning literature has focused primarily on code reuse, a truly effective library must also prioritize ease of use and maintainability, as well as be robust and adaptable to future extensions.

\textbf{Modularity and Code Quality}: %\todo{written code may not be of high quality, lack modularity}
LLMs are trained and optimized primarily for code correctness with insufficient focus on other aspects of code like quality and maintainability. This is further exacerbated with large scale reinforcement learning being performed using test cases which can lead to unintended consequences regarding code quality, as correct but poor quality code is still often given a high reward.
Empirically, it has been observed that LLM written solutions are often more complex than human-written counterparts. 
For example, \citet{jain2024r2e} identified that LLMs prefer to repeat existing code instead of making use of abstractions in the existing code. 
One aspect of code quality is modularity, ensuring that code does not get duplicated too often. Here, \citet{berlot2024library} identified that library or tool reuse is non-trivial for LLMs in coding and formal math domains. 

\textbf{Effective Tool Usage}: Software engineering has witnessed the development of various open and proprietary tooling support for programming, debugging, analysis, and code management over time.
For example, program analysis tools provide static and dynamic assurances on code correctness. 
Print statements and debuggers are used for dynamically analyzing and debugging programs at a fine-grained level.
Beyond programming, such tools are richly integrated into the entire software development lifecycle, e.g., code navigation or search, reviewing code, CI testing.

While many efforts combine LLMs and agents with tools, they do not achieve fully dynamic and effective software engineering tool usage. This involves an AI system seamlessly and proactively integrating appropriate tools depending on the task at hand. There are a few challenges towards achieving this goal. First, the AI system must identify which tools could potentially be useful for the task at hand. Second, the system then needs to decide when to invoke the tool. A complex debugging task might require the use of \texttt{pdb} or \texttt{gdb} to track intermediate program states, while looking at input-output pairs may be sufficient for simple debugging tasks. Third, the agent then must figure out how to invoke the tool. If the agent knows that a certain function in a program has an error, it may wish to walk through only that function instead of the entire code from start to finish. Finally, the agent needs to incorporate the output provided by the tool in order to inform its next steps, e.g. edit the code if a bug was uncovered or run the tool again otherwise.


\subsection{Interaction} \label{subsec:main-interaction}

While AI coding systems are increasingly more powerful, the majority of them are still at a low to medium autonomy level, serving as engineer assistance rather than achieving high or full automation. We identify a few key challenges of today's AI coding systems that prevent these systems from working with humans effectively at higher levels of autonomy.

\textbf{Vague Specifications and User Misalignment}: When using code LLMs or coding agents, we typically prompt them with a natural language specification. This can include a natural language description of the desired code, input-output examples, relevant code snippets, and other functional requirements. However, there is a gap in the level of abstraction between English and code, leading to incomplete or ambiguous specifications. This issue becomes more pronounced in longer programs, where the number of ambiguous decision points increases, and choices traditionally made by humans are instead implicitly embedded in the LLM’s generated code. Consequently, users often experience misalignment due to vague specifications. While many code LLMs support multi-turn interactions, it remains inherently challenging for users to articulate their thought processes into follow-up natural language instructions.

\textit{Inherent trade-offs in software development}: Designing large software systems always surfaces trade-offs between various desiderata such as readability, scalability, performance, maintainability, reliability, security, etc. These trade-offs are often context-dependent. A long-term and rapidly moving project may be willing to trade off some performance to have simplicity and readability. Performance-critical applications may completely sacrifice readability to eke out every millisecond of speed 
(such as using \href{https://graphics.stanford.edu/~seander/bithacks.html}{bit-twiddling hacks}).
Finding the sweet spot among these trade-offs can often involve extensive prototyping and benchmarking to understand the performance characteristics of different approaches. However, user specifications in the initial prompt rarely include details about these trade-offs, nor do models often take them into account.

% \textit{Implicit constraints}: Aside from functional/semantic correctness, there are also often implicit constraints in writing code. For example, many companies such as \href{https://opensource.janestreet.com/standards/}{Jane Street} and \href{https://google.github.io/styleguide/}{Google} have style guides, and many GitHub repositories explicitly outline style elements that new code ought to follow. \cite{zou2019does} find that GitHub pull requests that are more consistent with the style of the existing code get merged faster. Additionally, corporations may enforces codes of conduct or \href{https://developers.google.com/checks/guide/code-compliance/overview}{compliance requirements} at the code level. Furthermore, codebases follow common programming patterns or system design patterns that are implicitly specified by the way the current code is written. However, when using code LLMs, these constraints are often inferred incorrectly \citep{wang2024beyond}.

\textbf{Lack of Controllability}: When using AI coding systems, programmers often seek specific functionality, yet they lack reliable ways to steer LLMs toward generating precisely the desired code. Instead, they typically rely on a trial-and-error approach, repeatedly sampling outputs or providing feedback until the AI produces an acceptable solution. Consequently, significant human effort is still required to review and modify the code to ensure it meets the intended functionality \citep{weisz2024examining}.

A way to improve controllability is for AI coding systems to recognize when human input is needed and communicate effectively—-yet this remains the top-reported challenge in human-agent collaboration~\citep{shao2024collaborative}. LLMs rarely defer to humans for clarification, while developers often ask questions to clarify the description of a task provided by their peers. In addition, based on its knowledge of existing software, AI should incorporate implicit priors from a specification while keeping the user in the loop. 

% \textbf{Restricted Human-AI Interface}: Existing interfaces for code LLMs primarily manifest as intelligence features embedded within integrated development environments (IDEs). However, the ubiquitous ``Tab'' interaction paradigm prevalent in intelligent IDEs may prove inadequate when AI systems transition from completing developer-scaffolded functions to authoring the majority of the codebase autonomously. Current interfaces for coding agents, such as \href{https://devin.ai/}{Devin}, typically stream raw actions without adequate context or explanation. Given that these agents can execute numerous actions rapidly, developers face significant challenges in effectively monitoring the process, implementing timely interventions, or reasserting control when necessary. This lack of transparency can also undermine trust in AI-generated code~\citep{10.1145/3630106.3658984}.

\subsection{Measurement}

\textbf{Task Diversity and Capability Isolation}: 
Current coding evaluations primarily focus on the code generation task, with little attention towards the tasks discussed in Section~\ref{sec:main-tasks-milestones}.
As more agent-based approaches are introduced for software engineering (e.g. pairing a code generation model with a debugging model), these engineering-related capabilities beyond just code generation will be important towards designing a maximally performant system. 
Solely relying on end-to-end coding evaluations that focus on the overall correctness of a codebase makes it difficult to precisely measure progress and learn from the failure modes on individual tasks.


\textbf{Contamination}: 
Data contamination is a serious issue that, if not taken into account, can affect the soundness of various conclusions drawn from a set of benchmark results. In coding, the performance of LLMs on competitive programming \citep{xu2024benchmark, jain2024livecodebenchholisticcontaminationfree} and SWE-Bench \citep{aleithan2024swe} tasks has been shown to degrade over time, indicating the possibility of older problems being contaminated due to public exposure on the internet. For simpler function-level HumanEval style problems, \citet{matton2024leakage} suggest three potential causes of contamination: direct data leakage (benchmarks are on GitHub), synthetic data leakage (there are only a limited number of interview problems), and overfitting to test sets (benchmark hacking). In addition, for code, contamination can be hard to detect, as semantically equivalent code that is syntactically distinct could be thought of as contamination \citep{riddell2024quantifying}. A recent benchmark, the Konwinski Prize\footnote{\url{https://www.kprize.ai/}}, is a promising way to fairly evaluate SoTA LLM models by only using new GitHub issues.

\textbf{Construct Validity}:
Construct validity refers to how closely a measurement reflects the underlying concept. 
Given the implications of rapid performance improvement in the AI for the code domain, it is essential to have high-construct validity benchmarks evaluating how well programming agents can perform.
While benchmarks like SWE-Bench come close, user experiences do not currently match rapid performance gains obtained from them. 
This is partially because many desiderata in software engineering cannot be described cleanly via automated unit testing. Things like multi-turn code generation, designing an UI, and writing clean and idiomatic code are all difficult to quantitatively measure with precision. Designing reliable proxy metrics for these desired goals remains a challenge. 
